\section{Computational Intelligence nel mondo dei dati}
La \textbf{Computational Intelligence (CI)}, o intelligenza computazionale, è un 
ramo dell'Intelligenza Artificiale che si concentra sullo sviluppo di metodi 
computazionali in grado di apprendere, adattarsi, gestire l'incertezza e funzionare 
in modo efficace in ambienti molto complessi.
Questi metodi ottengono enormi successi traendo ispirazione diretta dall'intelligenza 
biologica e dal mondo naturale, creando quelli che vengono definiti 
"sistemi intelligenti".\\

Questi algoritmi intelligenti includono:
\begin{itemize}
    \item reti neurali artificiali
    \item calcolo evolutivo
    \item intelligenza dello sciame
    \item sistemi immunitari artificiali
    \item sistemi fuzzy
\end{itemize}

La fonte primaria ispiratrice della computational intelligence. Nasce come proposta
da un iraniano Zadeh, che propone la logica fuzzy, in cui diceva ottimi risultati potevano
essere ottenuti da singole applicazioni.

\subsection{Evoluzione della CI}
Nel mondo dei dati, la CI rappresenta un pilastro fondamentale per quattro motivi principali:
\begin{enumerate}
    \item \textbf{Gestione dell'incertezza}: Fornisce gli strumenti per elaborare 
    dati ambigui, rumorosi o incompleti, permettendo di prendere decisioni e fare 
    previsioni anche quando le informazioni non sono perfette.
    \item \textbf{Apprendimento e adattamento}: Le tecniche di CI estraggono 
    informazioni direttamente dai dati per adattare i propri parametri. 
    Questo è cruciale in un mondo dinamico dove i dati cambiano rapidamente nel tempo
    \item \textbf{Complessità}: Riesce a gestire e risolvere problemi altamente 
    complessi che risultano intrattabili con i metodi analitici e ingegneristici 
    tradizionali.
    \item \textbf{Automazione}: Automatizza processi e decisioni complesse, 
    liberando l'essere umano e permettendogli di concentrarsi su attività 
    più creative e strategiche.
\end{enumerate}

\subsection{Modelli Intelligenti vs. Modelli Algoritmici}
La vera rivoluzione della CI nel campo dei dati è il passaggio dai modelli algoritmici
a quelli intelligenti. 

\textbf{Modelli Algoritmici}\\
Questi modelli rappresentano l'approccio computazionale classico:
\begin{itemize}
    \item \textbf{Logica e Sequenze}: Seguono una sequenza logica di istruzioni predefinite. 
    Si basano su algoritmi rigorosi, che definiscono passo dopo passo le operazioni necessarie 
    per risolvere un problema.
    \item \textbf{Determinismo}: Sono strettamente deterministici. 
    Questo significa che, operando nelle medesime condizioni, produrranno sempre 
    risultati identici a fronte di uno stesso input.
    \item \textbf{Limitati alla Programmazione}: Sono vincolati in modo stringente al codice scritto dal programmatore. Per eseguire un determinato compito, richiedono che la sequenza di istruzioni sia inserita esplicitamente e dettagliatamente all'interno del programma.
\end{itemize}

\textbf{Modelli Intelligenti}
Questi modelli, tipici della Computational Intelligence, 
superano i limiti della programmazione rigida:
\begin{itemize}
    \item \textbf{Apprendimento e Adattamento}: La loro caratteristica principale è 
    la capacità di imparare dai dati e di adattarsi alle nuove informazioni, 
    il tutto senza essere esplicitamente programmati per farlo.
    \item \textbf{Gestione dell'Incertezza}: Mentre i modelli algoritmici richiedono 
    dati precisi, i modelli intelligenti possono operare in ambienti complessi 
    gestendo l'incertezza. Sono in grado di prendere decisioni e fare previsioni 
    anche in presenza di dati rumorosi o incompleti.
    \item \textbf{Capacità Cognitive}: Riescono ad emulare alcune delle capacità cognitive umane, 
    come l'adattamento alle variazioni dell'ambiente, l'apprendimento continuo e il ragionamento.
\end{itemize}

\subsection{Apprendimento e Adattamento ai Dati}
Nell'ambito della Computational Intelligence (CI), l'apprendimento e l'adattamento 
sono concetti strettamente legati che traggono una forte ispirazione dai processi 
del mondo biologico e naturale.

\paragraph{Apprendimento (Inspirazione Biologica e Artificiale)} 
Nel mondo biologico, gli organismi viventi imparano continuamente dall'esperienza, 
come ad esempio un animale che impara a evitare un determinato cibo dopo averne 
provato gli effetti nocivi. Questo principio è stato traslato nell'Intelligenza Artificiale, 
trovando la sua massima espressione nelle reti neurali artificiali. 
In queste architetture, i "neuroni" (le unità di elaborazione) sono in grado di 
modificare la forza delle proprie connessioni in base ai dati di input che ricevono, 
consentendo alla rete di accumulare esperienza e "apprendere" nel tempo.
\paragraph{Adattamento e Adattamento Dinamico} 
In biologia, l'adattamento è il processo evolutivo attraverso cui una specie 
diventa progressivamente più idonea a sopravvivere nel suo ambiente 
(come gli uccelli che sviluppano becchi più adatti a raccogliere cibo difficile da 
raggiungere). Nei sistemi di CI, l'adattamento viene implementato tramite algoritmi 
che modificano il comportamento del sistema in base all'ambiente o ai nuovi dati in 
ingresso, migliorandone le prestazioni nel tempo.\\

Un concetto fondamentale in questo campo è l'\textbf{adattamento dinamico}, 
ovvero la capacità di un sistema di modificare i suoi parametri in tempo reale 
mentre è in esecuzione in un ambiente mutevole. 

Questo è un enorme vantaggio perché permette al sistema di adattarsi istantaneamente 
senza aver bisogno di essere fermato per un "re-training" (ri-addestramento) offline. 
I sistemi dotati di processi adattivi riducono drasticamente il tempo necessario per 
giungere a una soluzione rispetto ai vecchi metodi che dovevano calcolare o esplorare 
gran parte delle possibilità.

\begin{itemize}
    \item \textbf{Apprendimento Supervisionato}: L'algoritmo viene addestrato fornendogli 
    un set di dati in cui le risposte corrette (output desiderati) sono già note. 
    Il sistema si adatta calcolando l'errore rispetto all'output corretto e applicando 
    le necessarie correzioni. È molto utilizzato per problemi di classificazione e 
    regressione.
    \item \textbf{Apprendimento Non Supervisionato}: Al sistema vengono forniti dati 
    "grezzi", privi di etichette o output desiderati. L'algoritmo deve adattarsi 
    esplorando i dati senza una guida esterna, con lo scopo di identificare 
    autonomamente pattern, regole o strutture nascoste (utilizzato molto nel 
    clustering o nella riduzione dimensionale).
    \item \textbf{Apprendimento per Rinforzo}: Il sistema impara agendo all'interno 
    di un ambiente dinamico. L'algoritmo si adatta in base alle conseguenze delle 
    proprie azioni, ricevendo un feedback sotto forma di ricompensa (se l'azione è 
    giusta) o penalità (se è sbagliata).
\end{itemize}

\subsection{L'importanza della Generalizzazione}
Un concetto cardine quando si applica la CI ai dati è la \textbf{generalizzazione}. 
Un modello non deve limitarsi a memorizzare i dati di addestramento 
(fenomeno noto come overfitting), ma deve riuscire a estrarre regole e pattern 
profondi per applicarli con successo a nuovi dati mai visti prima. 
La generalizzazione è cruciale per l'efficienza del sistema, per la sua capacità 
di affrontare la variabilità e il rumore del mondo reale, e per sapersi adattare 
a contesti differenti. L'abilità di generalizzare viene solitamente misurata 
testando il modello su un set di dati separato da quello di addestramento.\\

\section{Soft Computing e Calcolo Bio-Ispirato}
Il \textbf{Soft Computing} è un campo strettamente legato all'Intelligenza 
Computazionale e definisce una forma di \textbf{computazione bio-ispirata}. 
Bio-ispirata principalmente per algpritmi genetici e reti neurali. 
Nella logica non c'è sempre una risposta univoca, per esempio colore dei capelli, sono
bianchi ma ci sono anche bianchi sporchi, bianchi con riflessi, bianchi con sfumature, ecc.

Si basa su algoritmi euristici che traggono la loro ispirazione direttamente dal mondo 
naturale, seguendo un approccio che viene anche definito "ottimizzazione naturale" o 
degli "algoritmi bio-ispirati".\\

L'obiettivo di questi algoritmi è replicare o adattarsi ai processi osservati in 
natura, imitando il comportamento degli organismi o dei sistemi biologici. 
Questo tipo di computazione si rivela particolarmente efficace per risolvere problemi 
molto complessi o per effettuare ottimizzazioni in contesti in cui le soluzioni devono
sapersi adattare a cambiamenti dinamici, esattamente come le specie in natura si 
adattano al proprio ambiente.\\

I principali metodi e algoritmi che rientrano sotto l'ombrello del Soft Computing includono:

\begin{itemize}
    \item \textbf{Reti Neurali Artificiali}: Ispirate alla struttura e al funzionamento del cervello umano,
    sono composte da unità di elaborazione chiamate "neuroni", collegate tra
    loro in schemi complessi, che collaborano per apprendere dai dati e fare 
    previsioni o prendere decisioni.
    \item \textbf{Algoritmi Genetici}: Basati sulla teoria dell'evoluzione di Darwin, permettono 
    di affrontare problemi di ottimizzazione di altissima complessità che 
    richiederebbero tempi impraticabili con gli algoritmi tradizionali.
    Utilizzano concetti come la selezione naturale del "più adatto", il 
    crossover (incrocio genetico) e la mutazione per generare una popolazione di 
    possibili soluzioni e farla evolvere progressivamente nel tempo.
    \item \textbf{Logica Fuzzy}: È un'estensione della logica classica che permette 
    di superare la rigidità degli insiemi binari (dove un valore è strettamente "vero" 
    o "falso") introducendo i "gradi di verità".
    La logica fuzzy permette di trattare concetti sfumati, imprecisi o vaghi (come
    "quasi vero" o "parzialmente falso"), consentendo ai sistemi di modellare molto meglio 
    l'ambiguità e la complessità del mondo reale.
    \item \textbf{Algoritmi euristici ispirati a comportamenti collettivi (Swarm Intelligence)}: 
    Sebbene raggruppati tra gli algoritmi naturali, spiccano modelli come la Particle 
    Swarm Optimization (PSO) e l'Ant Colony Optimization (ACO).
    La PSO replica il comportamento sociale degli stormi di uccelli o dei banchi 
    di pesci, i cui membri comunicano e condividono le proprie "scoperte" per guidare 
    l'intero gruppo verso l'obiettivo migliore. L'ACO, invece, si ispira al 
    comportamento delle formiche, che per risolvere problemi combinatori lasciano 
    tracce chimiche (feromoni) sul percorso per influenzare le decisioni del resto 
    della colonia. Sono utilizzati per risolvere problemi di ottimizzazione dei percorsi,
\end{itemize}

\subsection{Algoritmi bio-ispirati}
\paragraph{Problemi di ottimizzazione}
Mettiamo che dobbiamo affrontare un problema di ottimizzazione, ovvero trovare 
quei parametri che in uno spazio di soluzioni mi minimizzano o massimizzano una funzione obiettivo.
Qualora scrivo la ricerca del funzione la ricerca del funzionale per massimizzarlo, quel funzionale
si chiama fitness function, e il processo di ottimizzazione è chiamato processo di evoluzione.
Se devo minimizzare la chiamerò funzione costo.\\

Massimizzare questa funzione significa trovare delle caratteristiche indicative del problema, 
Determinare quali sono i parametri che mi permettono di massimizzare la funzione obiettivo,
e quindi trovare la soluzione migliore al problema che sto affrontando.\\

Si instaura un processo iterativo in cui, a ogni iterazione, vengono generate nuove 
soluzioni (o "individui") che vengono valutate in base alla loro "fitness" 
(cioè quanto bene soddisfano la funzione obiettivo). Le soluzioni migliori vengono 
selezionate per "riprodursi" e generare nuove soluzioni, che a loro volta vengono 
valutate e selezionate, e così via, fino a quando non si raggiunge una soluzione 
soddisfacente o si esaurisce il tempo disponibile.\\

Mettiamo che dobbiamo cercare un minimo globale di una funzione complessa, 
con molti minimi locali. Un algoritmo genetico, grazie alla sua capacità di esplorare 
ampi spazi di soluzioni e di evitare di rimanere intrappolato in minimi locali, 
può essere molto più efficace di un algoritmo tradizionale di ottimizzazione, 
che potrebbe facilmente rimanere bloccato in un minimo locale senza raggiungere 
il minimo globale.\\

Pure gli algoritmi bio-ispirati dipendono da dei parametri, come la dimensione della
popolazione, il tasso di mutazione, il tasso di crossover, ecc. Li posso rendere più
o meno esplorativi a seconda di come li setto, e questo è un aspetto molto importante
perché mi permette di bilanciare l'esplorazione (la capacità di cercare nuove soluzioni)
e lo sfruttamento (la capacità di migliorare le soluzioni già trovate).\\

A volte quando la mole di parametri è molto grande si cerca di privilegiare 
algoritmi meno precisi ma più veloci, per esempio la ricerca del gradiente.

\paragraph{Un'altra importante classificazione dei problemi di ottimizzazione}
Le euristiche (come gli Algoritmi Genetici, la Particle Swarn Optimization, ecc.)vengono 
suddivise in base a come bilanciano esplorazione e sfruttamento rispetto alle dimensioni
dello spazio delle soluzioni:

\begin{itemize}
    \item \textbf{Euristiche per VHDSS (Very High Dimension Solution Space)}: Quando 
    il problema è a dimensioni elevatissime e l'utente non ha alcuna idea di dove
    si trovi la soluzione, è necessario utilizzare algoritmi ad alta esplorazione. 
    In questi casi, il confronto premia il Calcolo Evolutivo (come gli Algoritmi
    Genetici) e in generale i metodi di Swarm Intelligence, capaci di mappare ampie
    porzioni di spazio sconosciuto.
    \item \textbf{Euristiche per HDSS (High Dimension Solution Space)}: Quando si ha
    solo un sospetto sulla zona in cui risiede la soluzione, le euristiche più
    performanti per l'esplorazione includono algoritmi come la Bacterial Chemotaxis, 
    il Simulated Annealing, la Tabu Search e l'uso della Logica Fuzzy.
    \item \textbf{Euristiche per LDSS (Low Dimension Solution Space)}: Quando si ha 
    già un'idea abbastanza precisa delle coordinate ottimali, la fase di esplorazione
    diventa secondaria rispetto a quella di sfruttamento (exploitation). In questa
    specifica fase di affinamento, il confronto premia algoritmi come la Particle
    Swarm Optimization (PSO) e la Flock of Starlings Optimization.
\end{itemize}

Explanability (spiegabilità) è un aspetto cruciale in questo contesto, soprattutto
quando si utilizzano modelli intelligenti che operano in ambienti complessi e dinamici.
La capacità di spiegare le decisioni e i processi di apprendimento di un sistema di
CI è fondamentale per garantire la fiducia degli utenti, facilitare la diagnosi di errori 
e migliorare la trasparenza complessiva del sistema.
Un sistema di CI con una buona explainability è in grado di fornire spiegazioni chiare
e comprensibili sulle ragioni alla base delle sue decisioni, sui dati che ha 
utilizzato per apprendere e sulle strategie che ha adottato per adattarsi a nuove 
situazioni.

Consideriamo l'euristica Particle Swarm Optimization (PSO) come esempio di algoritmo
bio-ispirato. Questi algoritmi hanno una ispirazione di intelligenza di sciame (ma anche 
sociale) per questo a lavorarci sopra c'erano anche sociologi, psicologi, biologi, ecc.
Ciascun individuo che farà parte dello sciame, verrà inizializzato nello spazio delle
soluzioni con una posizione e una velocità casuali. 
In base al punto in cui si trova ogni individuo misurerà la sua fitness, e in base
alla fitness degli altri individui, si muoverà verso la posizione migliore che ha 
trovato lui stesso e verso la posizione migliore trovata dagli altri individui.
In questo modo, lo sciame si muove collettivamente verso le aree dello spazio
delle soluzioni che presentano una fitness più elevata, esplorando e sfruttando
le informazioni condivise tra i membri dello sciame per trovare soluzioni ottimali o quasi ottimali.\\

Gli stimoli sono:
\begin{itemize}
    \item \textbf{La mia esperienza personale} e quindi l'informazione che ho accumulato fino a quel momento
    (la mia posizione nello spazio delle soluzioni e la mia fitness)
    \item \textbf{L'esperienza sociale}, ovvero l'informazione che condividono gli altri membri
    dello sciame
    \item \textbf{L'inerzia}, ovvero la tendenza a mantenere la propria velocità e direzione 
    attuale, che aiuta a bilanciare l'esplorazione e lo sfruttamento.
\end{itemize}
Sarà dentro un ciclo for questo:
$$v_k^j(t+1)=v_k^j(t)+\lambda^j(p_best^j_k-x_k^j(t))+\gamma^j(g_best^j-x_k^j(t))$$
k vuol dire che lo spazio è n dimensionale, j vuol dire che ci sono m individui nello
sciame, $v_k^j(t)$ è la velocità dell'individuo j nella dimensione k al tempo t,
$p_best$ è la posizione migliore che ha trovato lui stesso, $g_best$ è la posizione 
migliore trovata dagli altri membri dello sciame, $\lambda$ e $\gamma$ sono dei 
parametri che bilanciano l'importanza di questi due stimoli.\\
$$V_i^{k+1}=\omega V_i^k+c_1rand_1(...)x(pbest_i^k-s_i^k)+c_2rand_2(...)x(gbest^k-s_i^k)...$$
In questo modo, ogni individuo si muove verso la posizione migliore che ha trovato lui stesso
e verso la posizione migliore trovata dagli altri membri dello sciame, bilanciando
l'esplorazione e lo sfruttamento attraverso i parametri $\lambda$ e $\gamma$.

I parametri rand1 e rand2 sono dei numeri casuali che vengono generati ad ogni iterazione, e che
introducono una componente di casualità nel processo di aggiornamento della velocità,
che aiuta a evitare che lo sciame si blocchi in minimi locali e a favorire l'esplorazione 
di nuove aree dello spazio delle soluzioni.\\

\textbf{Diagramma di flusso che rappresenta l'Algoritmo Generale PSO}:

\section{Swarn intelligence as State Equations}
É una variazione importante perché si passa da un algoritmo stocastico a uno deterministico, 
rimandendo nell'ambito della swarm intelligence, ma con un approccio più matematico e formale.
Modo di approcciare in maniera supervisionato invece di lasciare la supervisione ai parametri stocastici.\\

Questa variazione é una variazione che la comunitá scientifica del calcolo evolutivo digerisce e non 
digerisce, perché la natura stessa del calcolo é sfumato, quindi qualsiasi introduzione di rigidità e 
determinismo viene vista come una forzatura. Questi algortimi spesso fanno riferimento a fenomeni
naturali ma non li trattano come dovrebbero essere trattati.\\

Siamo abituati a vederer l'espressione: 
$$v_k^j(t+1)=\omega v_k^j(t)+\lambda(p^j_{{best}_k}(t)-x_k^j(t))+\gamma(g_{best}(t)-x_k^j(t))+ \sum_{m=1}^N h_{km}v_m^j$$
In questo modo, ogni individuo si muove verso la posizione migliore che ha trovato lui stesso
e verso la posizione migliore trovata dagli altri membri dello sciame, bilanciando l'esplorazione e lo 
sfruttamento attraverso i parametri $\lambda$ e $\gamma$, e tenendo conto anche dell'influenza degli altri 
individui dello sciame attraverso il termine $\sum_{m=1}^N h_{km}v_m^j$.

Il parametro hgamma é il parametro sociale, $g_{\text{best}}$ (global best) é la posizione migliore trovata dall'intero sciame quindi non 
dovrebbe avere j all'apice.

Una volta aggiornata la velocitá aggiorniamo la posizione:
$$x_k^j(t+1)=x_k^j(t)+v_k^j(t+1)$$

Peró sto sommando una velocitá a una posizione, quindi non é propriamento corretto.

Per convertire questo algoritmo in un sistema dinamico, dobbiamo esprimere l'aggiornamento della velocità e 
della posizione in termini di equazioni differenziali o differenze discrete.
$$v_k^j(t+\Delta t)=\omega v_k^j(t)+\lambda(p_{{best}_k}(t)-x_k^j(t))+\gamma(g_{best}(t)-x_k^j(t))+ \sum_{m=1}^N h_{km}v_m^j$$
Nel momento in cui divido per $\Delta t$ ottengo:
$$\frac{v_k^j(t+\Delta t)-v_k^j(t)}{\Delta t}=\omega \frac{v_k^j(t)}{\Delta t}+\lambda\frac{p_{\text{best},k}(t)-x_k^j(t)}{\Delta t}+\gamma\frac{g_{\text{best}}(t)-x_k^j(t)}{\Delta t}+ \sum_{m=1}^N h_{km}\frac{v_m^j}{\Delta t}$$
Nel limite in cui $\Delta t$ tende a zero, otteniamo un sistema di equazioni differenziali ordinarie che descrivono 
l'evoluzione della velocità e della posizione degli individui dello sciame nel tempo.\\

La forzante ($F_k (t)$).

I vantaggi di usare le forme chiuse invece di un algoritmo stocastico sono molteplici:
\begin{itemize}
    \item \textbf{Analisi Matematica}: Le forme chiuse permettono di applicare strumenti di analisi matematica per studiare la stabilità, la convergenza e il comportamento dinamico del sistema, cosa che è molto più difficile da fare con algoritmi stocastici.
    \item \textbf{Controllo e Ottimizzazione}: Con un modello deterministico, è possibile progettare strategie di controllo e ottimizzazione più precise, poiché si ha una comprensione più chiara di come le variabili influenzano il sistema.
    \item \textbf{Prevedibilità}: I sistemi deterministici offrono una maggiore prevedibilità rispetto agli algoritmi stocastici, che possono produrre risultati diversi a ogni esecuzione a causa della loro natura casuale.
    \item \textbf{Efficienza Computazionale}: In alcuni casi, le forme chiuse possono essere più efficienti dal punto di vista computazionale, poiché non richiedono la generazione di numeri casuali o l'esecuzione di molte iterazioni per raggiungere una soluzione soddisfacente.
    \item \textbf{Applicabilità a Problemi Specifici}: In alcuni contesti, come il controllo di sistemi dinamici o l'ottimizzazione in tempo reale, un modello deterministico può essere più adatto e fornire risultati più affidabili rispetto a un algoritmo stocastico.
\end{itemize}

\section{Fuzzy logic}
Fuzzy significa rozzo, per distiunguerlo dalla logica booleana che é pulita e precisa. 
C'é il ragionamento approssimato, rappresneta un modo di riprodurre in modo artificiale i processi di 
ragionamento umano, che spesso non sono precisi e rigidi, ma piuttosto sfumati e imprecisi.\\

Non é la totale negazione della logica booleana, ma una sua estensione che permette di trattare concetti sfumati e imprecisi.
Alcuni concetti come caldo o freddo, alto o basso, sono intrinsecamente sfumati e non possono essere rappresentati in modo
preciso con la logica booleana.

Per esempio per il caldo o freddo passiamo da una rappresentazione numerica a una qualificazione linguistica.\\

Concetti fondamentali della logica fuzzy:
\begin{itemize}
    \item Modellare l'incertezza del linguaggio naturale
    \item Trattare il concetto di parziale verità
    \item Definire i valori compresi tra 0 e 1, ovvero tra il "completamente falso" e il "completamente vero"
\end{itemize}

Autore della fuzzy logic é Lotfi Zadeh, che ha introdotto il concetto di "fuzzy set" (insieme sfumato)
per rappresentare insiemi che non hanno confini netti, ma piuttosto gradi di appartenenza.\\

\subsection{Concetto di insieme}
Secondo la logica classica, un elemento o appartiene a un insieme o non appartiene, 
mentre nella logica fuzzy, un elemento può appartenere a un insieme con un certo grado di appartenenza, 
che è rappresentato da un numero compreso tra 0 e 1.\\

Potremmo dire che fuzzy é una finestrazione del booleano, ma in realtá un elemento puó appartenere in parte 
anche in due insiemi opposti, per esempio un individuo puó essere alto e basso allo stesso tempo, con gradi di appartenenza diversi.\\

L'insieme dato dalla logica classica si chiama "crisp set" (insieme nitido), mentre l'insieme dato dalla logica fuzzy 
si chiama "fuzzy set" (insieme sfumato).\\

\textbf{Diagramma di Venn} é l'operazione di intersezionei tra insiemi crisp.

Zadeh ha esteso il concetto di insieme a quello di insieme fuzzy, introducendo il concetto di "membership function" 
(funzione di appartenenza), che assegna a ogni elemento un grado di appartenenza all'insieme fuzzy, rappresentato 
da un numero compreso tra 0 e 1.\\
Se un elemento ha un grado di appartenenza di 0, significa che non appartiene affatto all'insieme,
mentre se ha un grado di appartenenza di 1, significa che appartiene completamente all'insieme, 
se é tra 0 e 1 appartiene parzialmente.\\

Le funzioni di Membership possono assumere diverse forme, come ad esempio funzioni triangolari, 
trapezoidali o gaussiane, a seconda del contesto e delle esigenze specifiche dell'applicazione.\\

Le funzioni di appartenenza possono essere ricondotte a 3 fattori principali:
\begin{itemize}
    \item altezza: ti da l'estremo superiore 
    \item core: ti da l'estremo inferiore 
    \item supporto: ti da l'intervallo di valori per cui la funzione di appartenenza è non nulla
\end{itemize}

Un numero fuzzy non rappresenta un valore preciso, ma piuttosto un intervallo di valori con gradi di appartenenza diversi.
Fuzzy singleton é un numero fuzzy che rappresenta un singolo valore con un grado di appartenenza di 1, 
mentre tutti gli altri valori hanno un grado di appartenenza di 0.\\

L'operazione di intersezione tra insiemi fuzzy é definita come il minimo dei gradi di appartenenza,
mentre l'operazione di unione é definita come il massimo dei gradi di appartenenza.\\

\subsection{Fuzzy systems}
I normali insiemi nella logica booleana vengono trasformati in insiemi fuzzy, che consistono nel suddividere
l'universo considerato in sottoinsieme definiti da funzioni di appartenenza.
La membership function assegna a ogni elemento un grado di appartenenza a ciascun sottoinsieme, 
che rappresenta la sua "appartenenza" a quel sottoinsieme.
Questo grado ha un valore compreso tra 0 e 1.\\

Un elemento puó avere piú di una membership function, e quindi appartenere a piú sottoinsiemi con gradi di appartenenza diversi.\\

\paragraph{Algoritmo di inferenza fuzzy}
\begin{itemize}
    \item Definire le variabili linguistiche e i loro insiemi fuzzy associati (initialization)
    \item Construire le membership function per ciascun insieme fuzzy (initialization)
    \item Costruire le regole basi (initialization): delle regole linguistiche if then che collegano le 
    variabili di input alle variabili di output, utilizzando operatori logici fuzzy (come AND, OR, NOT) per combinare le condizioni.
    \item Convertire i valori di input in gradi di appartenenza utilizzando le membership function (fuzzification)
    \item Valutazione delle regole secondo le regole base (inference)
    \item Combinare i risultati per ciascuna regola (inference)
    \item Convertire l'output fuzzy in un valore crisp utilizzando un metodo di defuzzificazione (defuzzification)
\end{itemize}

\paragraph{Fuzzy rules}
Sample fuzzy rules for air conditioning:\\
IF (temperature id cold OR too-cold) AND (target is warm) THEN command is heat.\\
IF (temperature is hot OR too-hot) AND (target is warm) THEN command is cool.\\
IF (temperature is warm) AND (target is warm) THEN command is no-change.\\
\paragraph{Operazioni Fuzzy}
OR:
MAX: $$\mu_{A \cup B}(x) = \max(\mu_A(x), \mu_B(x))$$
ASUM: $$\mu_{A \cup B}(x) = \mu_A(x) + \mu_B(x) - \mu_A(x) \cdot \mu_B(x)$$
BSUM: $$\mu_{A \cup B}(x) = \min(1, \mu_A(x) + \mu_B(x))$$
AND:
MIN: $$\mu_{A \cap B}(x) = \min(\mu_A(x), \mu_B(x))$$
PROD: $$\mu_{A \cap B}(x) = \mu_A(x) \cdot \mu_B(x)$$
BDIF: $$\mu_{A \cap B}(x) = \max(0, \mu_A(x) + \mu_B(x) - 1)$$

\paragraph{Accumulation methods}
Maximum: $$\mu_{\text{acc}}(x) = \max{\mu_A(x), \mu_B(x)}$$
Bounded sum: $$\mu_{\text{acc}}(x) = \min(1, \mu_A(x) + \mu_B(x))$$
Normalized sum: $$\mu_{\text{acc}}(x) = \frac{\mu_A(x) + \mu_B(x)}{\max{1, \max{\{\mu_A(x'), \mu_B(x')\}}}}$$

\paragraph{Esempio robot che segue una linea}

